\chapter{Unit 1: Foundations of Problem Solving and Design Thinking}


\section{Section A: Multiple Choice Practice Questions (60 Curated Problems)}

\noindent Practice and test your mastery across all core concepts of this unit. Each question reflects standard university examination patterns and Bloom's Taxonomy competencies.
\vspace{3mm}

\noindent\textbf{Q1. What is design thinking primarily used for?} \hfill [Tag: Introduction to Design Thinking as a structured approach to problem solving | Easy]
\begin{enumerate}[label=(\alph*),topsep=2pt,itemsep=1pt,partopsep=0pt,parsep=0pt]
    \item {Repeating existing solutions}
    \item {Avoiding user feedback}
    \item {Memorizing rules quickly}
    \item {Solving problems creatively}
\end{enumerate}
\vspace{2mm}

\noindent\textbf{Q2. Which activity is commonly associated with design thinking?} \hfill [Tag: Introduction to Design Thinking as a structured approach to problem solving | Easy]
\begin{enumerate}[label=(\alph*),topsep=2pt,itemsep=1pt,partopsep=0pt,parsep=0pt]
    \item {Choosing the first idea}
    \item {Ignoring different opinions}
    \item {Understanding user needs}
    \item {Avoiding experimentation}
\end{enumerate}
\vspace{2mm}

\noindent\textbf{Q3. What is the purpose of prototyping in design thinking?} \hfill [Tag: Introduction to Design Thinking as a structured approach to problem solving | Easy]
\begin{enumerate}[label=(\alph*),topsep=2pt,itemsep=1pt,partopsep=0pt,parsep=0pt]
    \item {Testing a possible solution}
    \item {Writing a final report}
    \item {Replacing user research}
    \item {Proving every idea is perfect}
\end{enumerate}
\vspace{2mm}

\noindent\textbf{Q4. Which sequence best represents a common design thinking process?} \hfill [Tag: Introduction to Design Thinking as a structured approach to problem solving | Easy]
\begin{enumerate}[label=(\alph*),topsep=2pt,itemsep=1pt,partopsep=0pt,parsep=0pt]
    \item {Test, purchase, define, copy, finish}
    \item {Measure, calculate, publish, grade, report}
    \item {Plan, advertise, sell, close, repeat}
    \item {Empathize, define, ideate, prototype, test}
\end{enumerate}
\vspace{2mm}

\noindent\textbf{Q5. What does ideation mean in design thinking?} \hfill [Tag: Introduction to Design Thinking as a structured approach to problem solving | Easy]
\begin{enumerate}[label=(\alph*),topsep=2pt,itemsep=1pt,partopsep=0pt,parsep=0pt]
    \item {Measuring finished results}
    \item {Collecting financial records}
    \item {Generating possible ideas}
    \item {Selecting a final product}
\end{enumerate}
\vspace{2mm}

\noindent\textbf{Q6. What is a simple problem?} \hfill [Tag: Understanding different types of problems | Easy]
\begin{enumerate}[label=(\alph*),topsep=2pt,itemsep=1pt,partopsep=0pt,parsep=0pt]
    \item {A problem with unknown participants}
    \item {A problem with a clear solution}
    \item {A problem that changes every minute}
    \item {A problem with competing goals}
\end{enumerate}
\vspace{2mm}

\noindent\textbf{Q7. Which problem has clearly defined goals and information?} \hfill [Tag: Understanding different types of problems | Easy]
\begin{enumerate}[label=(\alph*),topsep=2pt,itemsep=1pt,partopsep=0pt,parsep=0pt]
    \item {A well-structured problem}
    \item {An imaginary problem}
    \item {A random problem}
    \item {A hidden problem}
\end{enumerate}
\vspace{2mm}

\noindent\textbf{Q8. What is a complex problem?} \hfill [Tag: Understanding different types of problems | Easy]
\begin{enumerate}[label=(\alph*),topsep=2pt,itemsep=1pt,partopsep=0pt,parsep=0pt]
    \item {A problem with many connected factors}
    \item {A problem solved by one fixed rule}
    \item {A problem with no possible effects}
    \item {A problem requiring no judgment}
\end{enumerate}
\vspace{2mm}

\noindent\textbf{Q9. What is an ill-structured problem?} \hfill [Tag: Understanding different types of problems | Easy]
\begin{enumerate}[label=(\alph*),topsep=2pt,itemsep=1pt,partopsep=0pt,parsep=0pt]
    \item {A problem with unclear goals}
    \item {A problem with one known answer}
    \item {A problem with complete information}
    \item {A problem with fixed instructions}
\end{enumerate}
\vspace{2mm}

\noindent\textbf{Q10. Why should a problem be classified before solving it?} \hfill [Tag: Understanding different types of problems | Easy]
\begin{enumerate}[label=(\alph*),topsep=2pt,itemsep=1pt,partopsep=0pt,parsep=0pt]
    \item {To make every problem identical}
    \item {To avoid gathering information}
    \item {To guarantee one correct answer}
    \item {To choose a suitable approach}
\end{enumerate}
\vspace{2mm}

\noindent\textbf{Q11. What does problem framing mean?} \hfill [Tag: Problem framing | Easy]
\begin{enumerate}[label=(\alph*),topsep=2pt,itemsep=1pt,partopsep=0pt,parsep=0pt]
    \item {Removing all possible solutions}
    \item {Defining how a problem is viewed}
    \item {Writing a problem in capital letters}
    \item {Decorating a problem statement}
\end{enumerate}
\vspace{2mm}

\noindent\textbf{Q12. Which question can help frame a problem?} \hfill [Tag: Problem framing | Easy]
\begin{enumerate}[label=(\alph*),topsep=2pt,itemsep=1pt,partopsep=0pt,parsep=0pt]
    \item {Who is affected by it?}
    \item {When can research be avoided?}
    \item {Why should facts be removed?}
    \item {How can it be ignored?}
\end{enumerate}
\vspace{2mm}

\noindent\textbf{Q13. Which is the best example of a user-focused problem statement?} \hfill [Tag: Problem framing | Easy]
\begin{enumerate}[label=(\alph*),topsep=2pt,itemsep=1pt,partopsep=0pt,parsep=0pt]
    \item {How might we increase our profits?}
    \item {How might we use more technology?}
    \item {How might we finish the project?}
    \item {How might we help commuters find buses?}
\end{enumerate}
\vspace{2mm}

\noindent\textbf{Q14. What is a benefit of writing a clear problem statement?} \hfill [Tag: Problem framing | Easy]
\begin{enumerate}[label=(\alph*),topsep=2pt,itemsep=1pt,partopsep=0pt,parsep=0pt]
    \item {It limits all creative thinking}
    \item {It guarantees immediate success}
    \item {It guides the search for solutions}
    \item {It removes the need for research}
\end{enumerate}
\vspace{2mm}

\noindent\textbf{Q15. What should a good problem frame avoid?} \hfill [Tag: Problem framing | Easy]
\begin{enumerate}[label=(\alph*),topsep=2pt,itemsep=1pt,partopsep=0pt,parsep=0pt]
    \item {Describing the user need}
    \item {Assuming one solution too early}
    \item {Considering the problem context}
    \item {Identifying affected people}
\end{enumerate}
\vspace{2mm}

\noindent\textbf{Q16. What is the main focus of human-centered thinking?} \hfill [Tag: Human-centered thinking | Easy]
\begin{enumerate}[label=(\alph*),topsep=2pt,itemsep=1pt,partopsep=0pt,parsep=0pt]
    \item {Products' appearance alone}
    \item {Machines' speed and power}
    \item {People's needs and experiences}
    \item {Organizations' internal structure}
\end{enumerate}
\vspace{2mm}

\noindent\textbf{Q17. Which method can help understand users' experiences?} \hfill [Tag: Human-centered thinking | Easy]
\begin{enumerate}[label=(\alph*),topsep=2pt,itemsep=1pt,partopsep=0pt,parsep=0pt]
    \item {Copying a competitor}
    \item {Guessing without evidence}
    \item {Ignoring user comments}
    \item {Conducting user interviews}
\end{enumerate}
\vspace{2mm}

\noindent\textbf{Q18. What does empathy mean in design thinking?} \hfill [Tag: Human-centered thinking | Easy]
\begin{enumerate}[label=(\alph*),topsep=2pt,itemsep=1pt,partopsep=0pt,parsep=0pt]
    \item {Understanding another person's perspective}
    \item {Focusing only on business goals}
    \item {Agreeing with every personal opinion}
    \item {Designing without meeting users}
\end{enumerate}
\vspace{2mm}

\noindent\textbf{Q19. Why is user feedback important?} \hfill [Tag: Human-centered thinking | Easy]
\begin{enumerate}[label=(\alph*),topsep=2pt,itemsep=1pt,partopsep=0pt,parsep=0pt]
    \item {It eliminates the need for testing}
    \item {It makes every idea successful}
    \item {It replaces understanding the problem}
    \item {It helps improve a proposed solution}
\end{enumerate}
\vspace{2mm}

\noindent\textbf{Q20. Which action best shows human-centered thinking?} \hfill [Tag: Human-centered thinking | Easy]
\begin{enumerate}[label=(\alph*),topsep=2pt,itemsep=1pt,partopsep=0pt,parsep=0pt]
    \item {Observing how people use a service}
    \item {Avoiding users during development}
    \item {Choosing features without evidence}
    \item {Designing only from personal preference}
\end{enumerate}
\vspace{2mm}

\noindent\textbf{Q21. A team begins a project by observing commuters, interviewing them, and identifying frustrations before proposing solutions. Which design thinking principle is being applied?} \hfill [Tag: Introduction to Design Thinking as a structured approach to problem solving | Medium]
\begin{enumerate}[label=(\alph*),topsep=2pt,itemsep=1pt,partopsep=0pt,parsep=0pt]
    \item {Defining outcomes alone}
    \item {Estimating costs first}
    \item {Selecting technology first}
    \item {Starting with user needs}
\end{enumerate}
\vspace{2mm}

\noindent\textbf{Q22. A prototype receives negative feedback from users. What is the most appropriate design thinking response?} \hfill [Tag: Introduction to Design Thinking as a structured approach to problem solving | Medium]
\begin{enumerate}[label=(\alph*),topsep=2pt,itemsep=1pt,partopsep=0pt,parsep=0pt]
    \item {Defend the original idea}
    \item {Revise the prototype}
    \item {Ignore the feedback}
    \item {End the project immediately}
\end{enumerate}
\vspace{2mm}

\noindent\textbf{Q23. Which sequence best represents a common design thinking process?} \hfill [Tag: Introduction to Design Thinking as a structured approach to problem solving | Medium]
\begin{enumerate}[label=(\alph*),topsep=2pt,itemsep=1pt,partopsep=0pt,parsep=0pt]
    \item {Empathize, define, ideate, prototype, test}
    \item {Define, budget, launch, advertise, review}
    \item {Plan, approve, produce, distribute, measure}
    \item {Research, purchase, build, sell, expand}
\end{enumerate}
\vspace{2mm}

\noindent\textbf{Q24. Why does design thinking encourage teams to generate several possible solutions before choosing one?} \hfill [Tag: Introduction to Design Thinking as a structured approach to problem solving | Medium]
\begin{enumerate}[label=(\alph*),topsep=2pt,itemsep=1pt,partopsep=0pt,parsep=0pt]
    \item {To avoid user research}
    \item {To increase creative alternatives}
    \item {To delay decision-making}
    \item {To guarantee technical success}
\end{enumerate}
\vspace{2mm}

\noindent\textbf{Q25. A team has limited time and creates a simple paper model of a new clinic layout. What is the main purpose of this prototype?} \hfill [Tag: Introduction to Design Thinking as a structured approach to problem solving | Medium]
\begin{enumerate}[label=(\alph*),topsep=2pt,itemsep=1pt,partopsep=0pt,parsep=0pt]
    \item {To test an idea cheaply}
    \item {To avoid involving users}
    \item {To replace final construction}
    \item {To prove the design is perfect}
\end{enumerate}
\vspace{2mm}

\noindent\textbf{Q26. A company knows that customer complaints are increasing but cannot identify the underlying cause. What type of problem is this primarily?} \hfill [Tag: Understanding different types of problems | Medium]
\begin{enumerate}[label=(\alph*),topsep=2pt,itemsep=1pt,partopsep=0pt,parsep=0pt]
    \item {An ill-defined problem}
    \item {A well-defined problem}
    \item {A fixed procedure}
    \item {A routine calculation}
\end{enumerate}
\vspace{2mm}

\noindent\textbf{Q27. Which situation is most likely a well-defined problem?} \hfill [Tag: Understanding different types of problems | Medium]
\begin{enumerate}[label=(\alph*),topsep=2pt,itemsep=1pt,partopsep=0pt,parsep=0pt]
    \item {Calculating a required inventory total}
    \item {Improving neighborhood trust}
    \item {Reducing loneliness among students}
    \item {Redesigning a city park experience}
\end{enumerate}
\vspace{2mm}

\noindent\textbf{Q28. A public transportation issue involves affordability, environmental impact, convenience, and conflicting stakeholder priorities. How is this problem best classified?} \hfill [Tag: Understanding different types of problems | Medium]
\begin{enumerate}[label=(\alph*),topsep=2pt,itemsep=1pt,partopsep=0pt,parsep=0pt]
    \item {Complex and interconnected}
    \item {Fully predetermined}
    \item {Purely mathematical}
    \item {Simple and repetitive}
\end{enumerate}
\vspace{2mm}

\noindent\textbf{Q29. A team assumes that slow service is caused by too few employees, but observation shows that unclear procedures create most delays. What error did the team make?} \hfill [Tag: Understanding different types of problems | Medium]
\begin{enumerate}[label=(\alph*),topsep=2pt,itemsep=1pt,partopsep=0pt,parsep=0pt]
    \item {They confused symptoms with causes}
    \item {They prototyped too early}
    \item {They measured the wrong units}
    \item {They selected too many users}
\end{enumerate}
\vspace{2mm}

\noindent\textbf{Q30. Why can a technically correct solution fail when addressing a complex social problem?} \hfill [Tag: Understanding different types of problems | Medium]
\begin{enumerate}[label=(\alph*),topsep=2pt,itemsep=1pt,partopsep=0pt,parsep=0pt]
    \item {Technical methods are always unsuitable}
    \item {People respond in predictable ways}
    \item {The solution may ignore human behavior}
    \item {Complex problems lack any constraints}
\end{enumerate}
\vspace{2mm}

\noindent\textbf{Q31. Which problem statement is most likely to support productive ideation?} \hfill [Tag: Problem framing | Medium]
\begin{enumerate}[label=(\alph*),topsep=2pt,itemsep=1pt,partopsep=0pt,parsep=0pt]
    \item {Night travel is impossible to improve.}
    \item {Why are commuters careless at night?}
    \item {Build a larger security system immediately.}
    \item {How might we help commuters feel safer at night?}
\end{enumerate}
\vspace{2mm}

\noindent\textbf{Q32. A team changes the statement "How can we install more signs?" to "How can we help visitors find rooms confidently?" What improvement has occurred?} \hfill [Tag: Problem framing | Medium]
\begin{enumerate}[label=(\alph*),topsep=2pt,itemsep=1pt,partopsep=0pt,parsep=0pt]
    \item {The frame replaces research with opinion}
    \item {The frame narrows ideas to signage}
    \item {The frame eliminates all constraints}
    \item {The frame focuses on a user outcome}
\end{enumerate}
\vspace{2mm}

\noindent\textbf{Q33. Which question would best test whether a problem frame is too narrow?} \hfill [Tag: Problem framing | Medium]
\begin{enumerate}[label=(\alph*),topsep=2pt,itemsep=1pt,partopsep=0pt,parsep=0pt]
    \item {Does it avoid describing the user?}
    \item {Does it permit several solution paths?}
    \item {Does it name one preferred product?}
    \item {Does it repeat the initial assumption?}
\end{enumerate}
\vspace{2mm}

\noindent\textbf{Q34. A school defines its challenge as "students need better grades." Which revision is more actionable and human-centered?} \hfill [Tag: Problem framing | Medium]
\begin{enumerate}[label=(\alph*),topsep=2pt,itemsep=1pt,partopsep=0pt,parsep=0pt]
    \item {How might we rank students more accurately?}
    \item {How might we help students practice difficult concepts regularly?}
    \item {How might we purchase the newest textbooks?}
    \item {How might we require students to study longer?}
\end{enumerate}
\vspace{2mm}

\noindent\textbf{Q35. Why should a team revisit its problem statement after conducting user research?} \hfill [Tag: Problem framing | Medium]
\begin{enumerate}[label=(\alph*),topsep=2pt,itemsep=1pt,partopsep=0pt,parsep=0pt]
    \item {Users always provide complete solutions}
    \item {Research makes framing unnecessary}
    \item {The original frame must remain unchanged}
    \item {Research may reveal different priorities}
\end{enumerate}
\vspace{2mm}

\noindent\textbf{Q36. A designer notices that older adults avoid a digital service because they fear making irreversible mistakes. Which response best demonstrates human-centered thinking?} \hfill [Tag: Human-centered thinking | Medium]
\begin{enumerate}[label=(\alph*),topsep=2pt,itemsep=1pt,partopsep=0pt,parsep=0pt]
    \item {Remove all instructions from the interface}
    \item {Increase the number of advanced settings}
    \item {Require users to complete training first}
    \item {Add clear undo and confirmation features}
\end{enumerate}
\vspace{2mm}

\noindent\textbf{Q37. Why is observing users in their actual environment often valuable?} \hfill [Tag: Human-centered thinking | Medium]
\begin{enumerate}[label=(\alph*),topsep=2pt,itemsep=1pt,partopsep=0pt,parsep=0pt]
    \item {It guarantees unbiased interpretation}
    \item {It reveals behaviors users may not report}
    \item {It proves the first design is correct}
    \item {It removes the need for interviews}
\end{enumerate}
\vspace{2mm}

\noindent\textbf{Q38. A team designs a service for people with limited internet access but tests it only with university students who have reliable broadband. What is the main concern?} \hfill [Tag: Human-centered thinking | Medium]
\begin{enumerate}[label=(\alph*),topsep=2pt,itemsep=1pt,partopsep=0pt,parsep=0pt]
    \item {The prototype is too inexpensive}
    \item {The service has too many stakeholders}
    \item {The participants do not represent target users}
    \item {The team used qualitative research}
\end{enumerate}
\vspace{2mm}

\noindent\textbf{Q39. During an interview, a participant says a form is easy to complete, but repeatedly asks for help while using it. What should the team conclude?} \hfill [Tag: Human-centered thinking | Medium]
\begin{enumerate}[label=(\alph*),topsep=2pt,itemsep=1pt,partopsep=0pt,parsep=0pt]
    \item {The form is successful without changes}
    \item {The participant's statement is always sufficient}
    \item {The observed behavior deserves investigation}
    \item {The participant should be excluded}
\end{enumerate}
\vspace{2mm}

\noindent\textbf{Q40. Which action best demonstrates empathy during the design process?} \hfill [Tag: Human-centered thinking | Medium]
\begin{enumerate}[label=(\alph*),topsep=2pt,itemsep=1pt,partopsep=0pt,parsep=0pt]
    \item {Choosing the cheapest feature first}
    \item {Measuring success only through sales}
    \item {Assuming designers share users' needs}
    \item {Understanding users' goals and constraints}
\end{enumerate}
\vspace{2mm}

\noindent\textbf{Q41. A team begins with a fixed technical solution, interviews only users who already favor it, and treats contradictory feedback as resistance. Which design-thinking failure is most fundamental?} \hfill [Tag: Introduction to Design Thinking as a structured approach to problem solving | Hard]
\begin{enumerate}[label=(\alph*),topsep=2pt,itemsep=1pt,partopsep=0pt,parsep=0pt]
    \item {Skipping prototype evaluation}
    \item {Using quantitative evidence too early}
    \item {Confusing ideation with implementation}
    \item {Prematurely converging on a solution}
\end{enumerate}
\vspace{2mm}

\noindent\textbf{Q42. A hospital wants to reduce missed appointments. Which sequence best reflects a design-thinking approach when the causes are initially unclear?} \hfill [Tag: Introduction to Design Thinking as a structured approach to problem solving | Hard]
\begin{enumerate}[label=(\alph*),topsep=2pt,itemsep=1pt,partopsep=0pt,parsep=0pt]
    \item {Benchmark competitors, automate reminders, then survey patients}
    \item {Define a target metric, select software, then interview staff}
    \item {Observe appointment journeys, synthesize patterns, prototype interventions, then test}
    \item {Model scheduling capacity, optimize resources, then validate the workflow}
\end{enumerate}
\vspace{2mm}

\noindent\textbf{Q43. A prototype receives enthusiastic feedback from six participants, but all six were recruited from the design team's professional network. What is the strongest conclusion?} \hfill [Tag: Introduction to Design Thinking as a structured approach to problem solving | Hard]
\begin{enumerate}[label=(\alph*),topsep=2pt,itemsep=1pt,partopsep=0pt,parsep=0pt]
    \item {The prototype has generated a promising but weak signal}
    \item {The prototype should proceed directly to implementation}
    \item {The prototype has demonstrated product-market fit}
    \item {The prototype is invalid because qualitative evidence is insufficient}
\end{enumerate}
\vspace{2mm}

\noindent\textbf{Q44. During ideation, an engineer rejects an unusual concept because it violates current technical constraints. What is the most appropriate response?} \hfill [Tag: Introduction to Design Thinking as a structured approach to problem solving | Hard]
\begin{enumerate}[label=(\alph*),topsep=2pt,itemsep=1pt,partopsep=0pt,parsep=0pt]
    \item {Record it separately and never revisit it}
    \item {Accept it as final because creativity has priority}
    \item {Defer feasibility judgment while preserving the concept}
    \item {Reject it immediately to preserve feasibility}
\end{enumerate}
\vspace{2mm}

\noindent\textbf{Q45. A team repeatedly alternates between interviews, reframing, prototyping, and testing rather than following a one-way sequence. Which interpretation is most accurate?} \hfill [Tag: Introduction to Design Thinking as a structured approach to problem solving | Hard]
\begin{enumerate}[label=(\alph*),topsep=2pt,itemsep=1pt,partopsep=0pt,parsep=0pt]
    \item {The team is avoiding the need for problem definition}
    \item {The team lacks a valid design process}
    \item {The team is applying iterative convergence and divergence}
    \item {The team is replacing evidence with intuition}
\end{enumerate}
\vspace{2mm}

\noindent\textbf{Q46. A city's traffic congestion persists despite expanding roads because new capacity encourages additional car use. This problem is best classified as:} \hfill [Tag: Understanding different types of problems | Hard]
\begin{enumerate}[label=(\alph*),topsep=2pt,itemsep=1pt,partopsep=0pt,parsep=0pt]
    \item {A technical problem caused by inadequate measurement}
    \item {A complicated problem with one optimal answer}
    \item {A complex problem with adaptive system behavior}
    \item {A simple problem with a delayed solution}
\end{enumerate}
\vspace{2mm}

\noindent\textbf{Q47. Which characteristic most strongly distinguishes a complex problem from a merely complicated one?} \hfill [Tag: Understanding different types of problems | Hard]
\begin{enumerate}[label=(\alph*),topsep=2pt,itemsep=1pt,partopsep=0pt,parsep=0pt]
    \item {It contains many interacting parts}
    \item {Its outcomes change as stakeholders respond}
    \item {It requires knowledge from several disciplines}
    \item {Its solution requires substantial resources}
\end{enumerate}
\vspace{2mm}

\noindent\textbf{Q48. A software defect can be isolated, reproduced, and fixed using a known engineering method. Which problem type and approach are most appropriate?} \hfill [Tag: Understanding different types of problems | Hard]
\begin{enumerate}[label=(\alph*),topsep=2pt,itemsep=1pt,partopsep=0pt,parsep=0pt]
    \item {Complicated; decompose the defect and apply expert analysis}
    \item {Wicked; negotiate a temporary definition with stakeholders}
    \item {Ill-structured; begin with unrestricted brainstorming}
    \item {Complex; probe the system through repeated social experiments}
\end{enumerate}
\vspace{2mm}

\noindent\textbf{Q49. A homelessness initiative has no universally accepted definition of success, and each proposed intervention changes stakeholder expectations. What is the key implication?} \hfill [Tag: Understanding different types of problems | Hard]
\begin{enumerate}[label=(\alph*),topsep=2pt,itemsep=1pt,partopsep=0pt,parsep=0pt]
    \item {A larger dataset will reveal the single correct solution}
    \item {The team should prioritize the fastest technically feasible option}
    \item {A detailed project plan should precede stakeholder engagement}
    \item {The team should treat the situation as a wicked problem}
\end{enumerate}
\vspace{2mm}

\noindent\textbf{Q50. A school reports low attendance as the problem, but interviews reveal transportation, safety, family-work demands, and student disengagement interact differently across groups. What should the team do first?} \hfill [Tag: Understanding different types of problems | Hard]
\begin{enumerate}[label=(\alph*),topsep=2pt,itemsep=1pt,partopsep=0pt,parsep=0pt]
    \item {Rank the causes by frequency and address the largest}
    \item {Increase attendance incentives before investigating further}
    \item {Map the interacting system and compare stakeholder perspectives}
    \item {Select one cause and define a narrowly scoped intervention}
\end{enumerate}
\vspace{2mm}

\noindent\textbf{Q51. Which reframing best converts the statement "How do we reduce customer-support costs?" into a human-centered design opportunity?} \hfill [Tag: Problem framing | Hard]
\begin{enumerate}[label=(\alph*),topsep=2pt,itemsep=1pt,partopsep=0pt,parsep=0pt]
    \item {How might we eliminate complex customer requests?}
    \item {How might we help customers resolve routine issues confidently?}
    \item {How might we require customers to use self-service tools?}
    \item {How might we minimize the number of support employees?}
\end{enumerate}
\vspace{2mm}

\noindent\textbf{Q52. A team frames a problem as "How can we build a faster check-in kiosk?" after observing that elderly visitors avoid existing kiosks. What is the primary flaw?} \hfill [Tag: Problem framing | Hard]
\begin{enumerate}[label=(\alph*),topsep=2pt,itemsep=1pt,partopsep=0pt,parsep=0pt]
    \item {It measures speed rather than total operating cost}
    \item {It assumes the kiosk should remain the central solution}
    \item {It excludes technical feasibility from the initial question}
    \item {It focuses too much on the organization's objectives}
\end{enumerate}
\vspace{2mm}

\noindent\textbf{Q53. A design team uses the frame "How might we make medication use effortless for everyone?" Which revision is strongest?} \hfill [Tag: Problem framing | Hard]
\begin{enumerate}[label=(\alph*),topsep=2pt,itemsep=1pt,partopsep=0pt,parsep=0pt]
    \item {How might we improve medication adherence for patients?}
    \item {How might we create a universal medication-management platform?}
    \item {How might we make medication reminders more frequent?}
    \item {How might we help adults with changing schedules take critical medication reliably?}
\end{enumerate}
\vspace{2mm}

\noindent\textbf{Q54. Two stakeholders disagree: managers define the issue as low employee productivity, while employees describe unnecessary workflow friction. Which framing practice best resolves the disagreement?} \hfill [Tag: Problem framing | Hard]
\begin{enumerate}[label=(\alph*),topsep=2pt,itemsep=1pt,partopsep=0pt,parsep=0pt]
    \item {Use the manager's framing because organizational authority determines scope}
    \item {Choose the framing supported by the larger stakeholder group}
    \item {Combine both labels into a broader statement about productivity}
    \item {Investigate underlying needs and construct a testable shared opportunity}
\end{enumerate}
\vspace{2mm}

\noindent\textbf{Q55. A team narrows its challenge from "improve public health" to "help shift workers access nutritious meals during overnight shifts." What important trade-off has the team made?} \hfill [Tag: Problem framing | Hard]
\begin{enumerate}[label=(\alph*),topsep=2pt,itemsep=1pt,partopsep=0pt,parsep=0pt]
    \item {It increased solution freedom by removing the user context}
    \item {It improved actionability while reducing universal coverage}
    \item {It converted a complex problem into a purely technical problem}
    \item {It eliminated the need to validate the problem with users}
\end{enumerate}
\vspace{2mm}

\noindent\textbf{Q56. Users request an additional dashboard feature, but observation shows they rarely use the existing dashboard because they distrust the data. What should a human-centered team prioritize?} \hfill [Tag: Human-centered thinking | Hard]
\begin{enumerate}[label=(\alph*),topsep=2pt,itemsep=1pt,partopsep=0pt,parsep=0pt]
    \item {Improve visual design before investigating trust}
    \item {Investigate the causes of distrust before selecting features}
    \item {Implement the requested feature to demonstrate responsiveness}
    \item {Ignore stated requests because observed behavior is always superior}
\end{enumerate}
\vspace{2mm}

\noindent\textbf{Q57. A voice-based service works well for most users but exposes private information in shared households. Which response best reflects human-centered thinking?} \hfill [Tag: Human-centered thinking | Hard]
\begin{enumerate}[label=(\alph*),topsep=2pt,itemsep=1pt,partopsep=0pt,parsep=0pt]
    \item {Require every household to create separate physical rooms}
    \item {Keep the design because average usability is high}
    \item {Study affected contexts and redesign disclosure controls with users}
    \item {Remove voice interaction because privacy risks are unavoidable}
\end{enumerate}
\vspace{2mm}

\noindent\textbf{Q58. A team invites users to a workshop but gives them only predefined solution choices. Why might this fail to be genuinely human-centered?} \hfill [Tag: Human-centered thinking | Hard]
\begin{enumerate}[label=(\alph*),topsep=2pt,itemsep=1pt,partopsep=0pt,parsep=0pt]
    \item {The format constrains users from revealing unmet needs}
    \item {Users cannot provide statistically representative data}
    \item {Predefined choices always produce technically inferior designs}
    \item {The workshop is too collaborative for objective evaluation}
\end{enumerate}
\vspace{2mm}

\noindent\textbf{Q59. A financial app improves average completion time but causes users with low literacy to abandon applications more often. Which evaluation principle should guide the next iteration?} \hfill [Tag: Human-centered thinking | Hard]
\begin{enumerate}[label=(\alph*),topsep=2pt,itemsep=1pt,partopsep=0pt,parsep=0pt]
    \item {Increase application speed further to compensate for comprehension barriers}
    \item {Examine distributional effects and redesign around excluded users}
    \item {Optimize the mean because aggregate performance is most reliable}
    \item {Discard the app because any subgroup disparity proves total failure}
\end{enumerate}
\vspace{2mm}

\noindent\textbf{Q60. A researcher asks participants to describe how they normally use a transit ticketing system, but participants give socially desirable answers. Which method would best complement the interviews?} \hfill [Tag: Human-centered thinking | Hard]
\begin{enumerate}[label=(\alph*),topsep=2pt,itemsep=1pt,partopsep=0pt,parsep=0pt]
    \item {Observe real ticket-purchasing behavior in context}
    \item {Use only the responses that match the researcher's hypothesis}
    \item {Ask the same questions with more formal wording}
    \item {Replace participants with transportation experts}
\end{enumerate}
\vspace{2mm}


\newpage
\section{Section B: Comprehensive Subjective Problem with Detailed Analytical Rubric}

\begin{subjectiveq}{1}{1}
\textbf{Case Background:} A municipal transport corporation notices a 42\% decline in bus ridership among senior citizens and persons with physical disabilities despite heavy fare subsidies. A technical consulting firm proposed simply buying low-floor electric buses with an estimated expenditure of \$12 million. However, preliminary commuter interviews revealed that anxiety stemmed primarily from lack of real-time accessible arrival information, intimidating boarding interactions, and confusing transfers.

\textbf{Required Design Deliverables:}
\begin{enumerate}[label=(\alph*)]
    \item Analyze why this scenario qualifies as an \textit{ill-structured wicked problem} rather than a \textit{well-structured engineering problem}.
    \item Map out an end-to-end intervention using the Stanford d.school 5-stage Design Thinking framework (Empathize, Define, Ideate, Prototype, Test).
    \item Formulate a human-centered Point-of-View (POV) statement and derive three actionable ``How Might We'' (HMW) questions.
\end{enumerate}
\end{subjectiveq}

\begin{solbox}
\textbf{1. Analytical Classification (Wicked vs. Well-Structured Problem):}\\
Under the criteria defined by Rittel and Webber (1973), urban transit accessibility is an ill-structured wicked problem because:
\begin{itemize}
    \item \textit{No definitive formulation:} The root issue cannot be decoupled from human psychology, civic infrastructure, and digital accessibility.
    \item \textit{Non-technical failure mode:} The consulting firm committed the classic fallacy of treating a socio-technical challenge as a tame mechanical procurement problem.
    \item \textit{No stopping rule:} Any solution alters commuter behavioral patterns, revealing secondary ripple effects (e.g., driver training, station shelter safety).
\end{itemize}

\textbf{2. Five-Stage Stanford d.school Design Thinking Architecture:}\\
\begin{itemize}
    \item \textbf{Stage 1: Empathize:} Deploy contextual inquiry, shadow senior citizens across complete journeys, conduct extreme-user interviews, and construct an Empathy Map (Says, Thinks, Does, Feels).
    \item \textbf{Stage 2: Define:} Synthesize qualitative findings to isolate core emotional friction points (fear of falling, disorientation during unannounced gate changes).
    \item \textbf{Stage 3: Ideate:} Conduct cross-functional Brainwriting 6-3-5 sessions with bus operators, geriatric specialists, and transit engineers to generate non-capital-intensive interventions.
    \item \textbf{Stage 4: Prototype:} Develop low-fidelity tactile boarding tokens, high-contrast audio-visual bus-stop beacons, and a simplified single-tap companion mobile interface.
    \item \textbf{Stage 5: Test:} Run a 14-day controlled pilot across two high-density suburban routes, tracking Task Completion Rate, Usability Anxiety Score, and Driver Assistance frequency.
\end{itemize}

\textbf{3. Human-Centered Problem Statement (POV) \& HMW Questions:}\\
\textit{POV Statement:} ``Elderly and mobility-impaired urban commuters need predictable, dignity-preserving transit guidance and low-stress boarding support because unexpected transit disruptions and physical rush induce severe panic and social isolation.''

\textit{Derived HMW Questions:}
\begin{enumerate}[label=HMW \arabic*:]
    \item How might we communicate bus arrival certainty in high-contrast physical and auditory formats without requiring smartphones?
    \item How might we redesign the physical bus dwell and boarding interaction so seniors feel calm rather than rushed by queues?
    \item How might we incentivize and train bus captains to act as proactive accessibility navigators during boarding?
\end{enumerate}
\end{solbox}

\begin{examinertip}
\textbf{Criteria for Full 10/10 Marks:}\\
1. \textbf{Rigorous Theoretical Framing (3 Marks):} Explicitly identify and define the formal model or framework applied (e.g., d.school stages, 5-Whys, MCDA, SCAMPER, Blueprinting, Minto SCQA). Avoid vague hand-waving.\\
2. \textbf{Contextual Application (4 Marks):} Directly connect every step to the specific facts, users, and constraints given in the case scenario.\\
3. \textbf{Actionable Deliverables \& Invariants (3 Marks):} Provide precise, formatted tables, matrices, or SMART criteria with clear quantitative bounds.
\end{examinertip}
